%【紙面設定】
%b4,横書き,2段
\documentclass[paper=b4j,landscape,twocolumn,fleqn]{jlreq}
\special{papersize=\the\paperwidth,\the\paperheight}
\setlength{\columnseprule}{0.4pt}
\pagestyle{empty}
\usepackage[margin=10truemm]{geometry}
\usepackage[dvipdfmx]{graphicx}
\usepackage{amsmath}
\ModifyHeading{section}{lines=1}

%【本文】
\begin{document}
\section{中点 解答}
点A(2,3),点B(-2,1),点C(4,-9),点D(-4,-1)
\begin{align*}
&1.\ (0,2)\\
&2.\ (3,-3)\\
&3.\ (-1,1)\\
&4.\ (1,-4)\\
&5.\ (0,-5)\\
&6.\ \left(\frac{4}{3},-\frac{5}{3}\right)\\
&7.\ \left(-\frac{2}{3},-3\right)\\
&8.\ \left(0,-\frac{3}{2}\right)
\end{align*}

\section{2点間の距離 解答}
\begin{align*}
&1.\ 2\sqrt{5}\\
&2.\ 2\sqrt{34}\\
&3.\ 8\sqrt{2}\\
&4.\ 7\\
&5.\ \sqrt{29}\\
&6.\ \frac{2\sqrt{65}}{3}\\
&7.\ \frac{2\sqrt{97}}{3}\\
&8.\ \frac{3}{2}
\end{align*}

\newpage
\section{2点間の移動 解答}
(1秒あたりの移動量)
\begin{align*}
&1.\ (-2,-1)\\
&2.\ \left(\frac{1}{2},-3\right)\\
&3.\ (-3,-2)\\
&4.\ \left(\frac{3}{2},-\frac{5}{2}\right)\\
&5.\ (-1,-1)\\
&6.\ (-2,2)\\
&7.\ (2,1)\\
&8.\ \left(-\frac{1}{2},3\right)
\end{align*}

\section{円の当たり判定1 解答}
\begin{align*}
&1.\ |AB|=\sqrt{45}=3\sqrt{5}>1+2=3\ \Rightarrow\ \text{当たらない}\\
&2.\ |AC|=\sqrt{2}<1+3=4\ \Rightarrow\ \text{当たる。必要移動量}=4-\sqrt{2}
\end{align*}

\section{円の当たり判定2 解答}
\begin{align*}
&3.\ \text{Aの必要総移動}=4\ \Rightarrow\ 1\text{秒あたり }(2,0)\\
&4.\ \text{Bの必要総移動}=4\ \Rightarrow\ 1\text{秒あたり }(1,0)
\end{align*}

\section{円の当たり判定の一般化 解答}
\begin{align*}
&1.\ \sqrt{(a-c)^2+(b-d)^2}\le r+R\\
&2.\ (a-c)^2+(b-d)^2\le (r+R)^2
\end{align*}
\end{document}
